\documentclass[a4paper,12pt]{article}
\usepackage[T2A]{fontenc}
\usepackage[utf8]{inputenc}
\usepackage[russian]{babel}
\usepackage{amsmath,amssymb,mathtools}
\usepackage{helvet}
\renewcommand{\familydefault}{\sfdefault}
\usepackage{icomma}
\usepackage[a4paper,margin=18mm,top=17mm,bottom=18mm]{geometry}
\usepackage{microtype}
\usepackage{tikz}
\usetikzlibrary{calc,arrows.meta,positioning,shapes.geometric}
\usepackage{pgfplots}
\pgfplotsset{compat=1.18}
\usepackage{tabularx,booktabs,longtable}
\usepackage{enumitem}
\usepackage[hidelinks]{hyperref}
\usepackage{parskip}
\usepackage{fancyhdr}
\usepackage{setspace}
\usepackage{xcolor}

\definecolor{accent}{HTML}{0F6B6D}
\definecolor{darkaccent}{HTML}{07494B}
\definecolor{textgray}{HTML}{2E3133}
\definecolor{softgray}{HTML}{6B7075}
\definecolor{gridgray}{HTML}{D9DDDF}
\definecolor{softaccent}{HTML}{E8F2F1}
\definecolor{warn}{HTML}{8B5A2B}

\color{textgray}
\onehalfspacing
\setlength{\parindent}{0pt}
\setlist[itemize]{leftmargin=*,itemsep=2pt,topsep=3pt}
\setlist[enumerate]{leftmargin=*,itemsep=4pt,topsep=4pt}
\pagestyle{fancy}
\fancyhf{}
\renewcommand{\headrulewidth}{0pt}
\fancyfoot[C]{\small\color{softgray}\thepage}

\newcommand{\sectionline}{\par\vspace{2pt}\noindent\color{accent}\rule{\linewidth}{0.6pt}\par\color{textgray}\vspace{4pt}}
\newcommand{\term}[1]{\textbf{\color{darkaccent}#1}}
\newcommand{\checkitem}{\raisebox{0.15ex}{\tikz\draw[accent,line width=0.7pt] (0,0) rectangle (0.18,0.18);}\hspace{0.45em}}

\tikzset{
  flowstart/.style={ellipse,draw=accent,line width=.9pt,fill=softaccent,minimum width=36mm,minimum height=9mm,align=center,font=\small},
  flowact/.style={rounded corners=2mm,draw=accent,line width=.8pt,fill=white,text width=104mm,minimum height=11mm,align=center,font=\small},
  flowcond/.style={diamond,aspect=2.4,draw=accent,line width=.8pt,fill=white,text width=55mm,align=center,font=\small,inner sep=1.5pt},
  flowarr/.style={-{Latex[length=2.4mm]},line width=.8pt,draw=darkaccent}
}

\begin{document}

\begin{titlepage}
\thispagestyle{empty}
\centering
\vspace*{22mm}
{\Large\color{softgray}Профильная математика · ЕГЭ\par}
\vspace{15mm}
{\fontsize{28}{34}\selectfont\bfseries\color{darkaccent}Чек-лист\par}
\vspace{5mm}
{\fontsize{20}{25}\selectfont\bfseries Графики показательной функции\\[2mm]и метод узловых точек\par}
\vspace{18mm}
{\large Даниил Кичук\par}
\vspace{5mm}
{\large 02.10.26\par}
\vfill
{\normalsize\color{softgray}Лёвин Артём Александрович\par}
{\small\color{softgray}эксклюзивно для Даниила Кичука\par}
\vspace{12mm}
\begin{tikzpicture}[remember picture,overlay]
  \draw[accent!55,line width=1.2pt]
  ([xshift=18mm,yshift=15mm]current page.south west)
  .. controls ([xshift=62mm,yshift=27mm]current page.south west) and ([xshift=-60mm,yshift=7mm]current page.south east) ..
  ([xshift=-18mm,yshift=18mm]current page.south east);
\end{tikzpicture}
\end{titlepage}

\section*{1. Что нужно уметь после занятия}
\sectionline
\begin{itemize}
  \item \checkitem находить на графике \term{узловые точки} с точно читаемыми координатами;
  \item \checkitem выбирать точки, которые упрощают вычисления и действительно определяют неизвестные параметры;
  \item \checkitem переводить точку $(x_0,y_0)$ графика в равенство $f(x_0)=y_0$;
  \item \checkitem восстанавливать параметры функций вида $f(x)=a^x+b$ и $f(x)=a^{x+b}$;
  \item \checkitem использовать степени $a^0=1$, $a^{-n}=\dfrac1{a^n}$ и дробные показатели;
  \item \checkitem после восстановления формулы находить $f(k)$ или решать уравнение $f(x)=c$;
  \item \checkitem проверять найденную функцию по исходным точкам графика.
\end{itemize}

\section*{2. Узловая точка: что это и зачем она нужна}
\sectionline
\term{Узловая точка} — точка графика, координаты которой можно считать точно, обычно это целые числа. Если точка $(x_0,y_0)$ лежит на графике функции, то
\[
  y_0=f(x_0).
\]
Именно это равенство превращает рисунок в уравнение для неизвестных параметров.

\textbf{Сколько точек брать?} Обычно столько, сколько неизвестных параметров. Для $f(x)=a^x+b$ неизвестны $a$ и $b$, поэтому нужны две \emph{независимые} точки. Если одна точка уничтожает неизвестный параметр и не даёт новой информации, её заменяют.

\textbf{Какие точки удобнее?} В первую очередь точки с $x=0$, $y=0$, $x=\pm1$ и вообще с небольшими по модулю координатами. Далёкая точка тоже допустима, но вычисления часто становятся тяжелее.

\section*{3. Главная связка: график $\longrightarrow$ уравнения}
\sectionline
Для точки $(x_0,y_0)$ всегда соблюдайте порядок:
\[
  \boxed{x=x_0,\qquad f(x)=y_0.}
\]
Запись $y$ на координатной плоскости и запись $f(x)$ в формуле обозначают одно и то же значение функции.

\textbf{Частая ошибка:} брать точку, которая визуально близка к графику, но не лежит на нём точно. Для восстановления параметров подходят только точки самого графика.

\section*{4. Базовый пример: $f(x)=a^x+b$}
\sectionline
Пусть график проходит через точки $A(0,-2)$ и $B(1,-1)$. Требуется найти $f(6)$.

\begin{center}
\begin{tikzpicture}
\begin{axis}[
  width=0.72\linewidth,height=6.6cm,
  axis lines=middle,
  xmin=-2,xmax=4.2,ymin=-4,ymax=7,
  xtick={-2,-1,0,1,2,3,4},
  ytick={-4,-3,-2,-1,0,1,2,3,4,5,6},
  grid=both,minor tick num=0,
  grid style={draw=gridgray,line width=.3pt},
  axis line style={draw=softgray},
  tick label style={font=\small},
  clip=false
]
\addplot[domain=-2:3.25,samples=220,line width=1.4pt,color=accent]{2^x-3};
\addplot[only marks,mark=*,mark size=2.5pt,color=darkaccent] coordinates {(0,-2) (1,-1)};
\node[anchor=south east,font=\small] at (axis cs:0,-2) {$A(0,-2)$};
\node[anchor=south west,font=\small] at (axis cs:1,-1) {$B(1,-1)$};
\end{axis}
\end{tikzpicture}
\end{center}

Подставляем координаты двух точек:
\[
\begin{cases}
-2=a^0+b,\\
-1=a^1+b.
\end{cases}
\qquad\Longrightarrow\qquad
\begin{cases}
-2=1+b,\\
-1=a+b.
\end{cases}
\]
Отсюда $b=-3$, затем $a=2$. Значит,
\[
  f(x)=2^x-3,
  \qquad
  f(6)=2^6-3=61.
\]

\textbf{Почему точка с $x=0$ особенно полезна?} Потому что при допустимом основании $a>0$ имеем $a^0=1$, и один параметр часто находится почти мгновенно.

\section*{5. Если степень сама содержит параметр: $f(x)=a^{x+b}$}
\sectionline
Пусть известны точки $(-3,1)$ и $(-1,2)$.
\[
\begin{cases}
1=a^{-3+b},\\
2=a^{-1+b}.
\end{cases}
\]
Так как для показательной функции $a>0$, $a\ne1$, равенство $1=a^0$ позволяет сразу получить
\[
-3+b=0\quad\Longrightarrow\quad b=3.
\]
Тогда
\[
2=a^{-1+3}=a^2
\quad\Longrightarrow\quad
 a=\sqrt2=2^{1/2}.
\]
Поэтому
\[
 f(x)=\left(2^{1/2}\right)^{x+3}=2^{\frac{x+3}{2}}.
\]
Например,
\[
 f(-7)=2^{\frac{-7+3}{2}}=2^{-2}=\frac14.
\]

\section*{6. Степени: короткие приёмы, которые экономят время}
\sectionline
\begin{itemize}
  \item $a^0=1$ — используйте сразу, как только появляется нулевая степень.
  \item $1=a^0$ — полезно в обратную сторону, когда нужно сравнить показатели при одинаковом основании.
  \item Если $a^m=c$ и $m\ne0$, то $a=c^{1/m}$, с учётом условия $a>0$ для основания показательной функции.
  \item Например, $a^{-2}=4 \Rightarrow a=4^{-1/2}=\dfrac12$.
  \item Если основание дробное, пишите скобки: $\left(\dfrac12\right)^{x-3}$. Без скобок степень может относиться только к знаменателю или к последнему символу.
\end{itemize}

\clearpage
\section*{Алгоритм: восстановление функции по графику}
\vspace{8mm}
\begin{center}
\begin{tikzpicture}[node distance=7mm]
\node[flowstart] (s) {Начало};
\node[flowact,below=of s] (f) {Запишите общий вид функции и отметьте неизвестные параметры};
\node[flowact,below=of f] (p) {Выберите узловые точки с точно читаемыми и удобными координатами};
\node[flowcond,below=9mm of p] (q) {Точек достаточно и они дают независимые уравнения?};
\node[flowact,below=10mm of q] (sub) {Для каждой точки $(x_0,y_0)$ запишите $y_0=f(x_0)$};
\node[flowact,below=of sub] (solve) {Упростите степени и найдите параметры};
\node[flowact,below=of solve] (build) {Соберите формулу функции};
\node[flowact,below=of build] (target) {Выполните требуемое: найдите $f(k)$ или решите $f(x)=c$};
\node[flowact,below=of target] (check) {Проверьте формулу подстановкой исходных точек};
\node[flowstart,below=of check] (e) {Готово};

\draw[flowarr] (s)--(f);
\draw[flowarr] (f)--(p);
\draw[flowarr] (p)--(q);
\draw[flowarr] (q)--node[right,font=\small]{да}(sub);
\draw[flowarr] (sub)--(solve);
\draw[flowarr] (solve)--(build);
\draw[flowarr] (build)--(target);
\draw[flowarr] (target)--(check);
\draw[flowarr] (check)--(e);
\draw[flowarr] (q.east) -- ++(22mm,0) node[above,font=\small,pos=.35]{нет} |- (p.east);
\end{tikzpicture}
\end{center}
\clearpage

\section*{7. Точки, которые «не сработали»}
\sectionline
Для функции $f(x)=a^x$ точка $(0,1)$ верна при любом допустимом $a$:
\[
  1=a^0=1.
\]
Она подтверждает, что точка лежит на графике, но \textbf{не определяет основание}. Нужно взять другую узловую точку, например $(-1,2)$:
\[
  2=a^{-1}\quad\Longrightarrow\quad a=\frac12.
\]

\textbf{Правило:} если после подстановки неизвестный параметр исчез и определить его стало невозможно, выберите другую точку. Это нормальная часть метода.

\section*{8. Контрольные признаки корректного решения}
\sectionline
Перед ответом быстро проверьте:
\begin{itemize}
  \item \checkitem все выбранные точки действительно лежат на графике;
  \item \checkitem координаты не перепутаны: сначала $x$, затем $y=f(x)$;
  \item \checkitem число независимых уравнений достаточно для числа неизвестных;
  \item \checkitem использованы условия $a>0$ и $a\ne1$ для показательной функции;
  \item \checkitem дробное основание взято в скобки;
  \item \checkitem отрицательная степень обработана как обратная величина;
  \item \checkitem найденная формула возвращает исходные точки;
  \item \checkitem выполнено именно то, что спросили: $f(k)$ или значение $x$.
\end{itemize}

\clearpage
\section*{9. Типичные ошибки}
\sectionline
\begin{tabularx}{\linewidth}{@{}p{0.31\linewidth}X@{}}
\toprule
\textbf{Ошибка} & \textbf{Как исправить}\\
\midrule
Взяли точку рядом с графиком & Используйте только точку, точно лежащую на графике.\\
Перепутали $x$ и $y$ & Для $(x_0,y_0)$ всегда пишите $f(x_0)=y_0$.\\
Взяли слишком неудобные точки & Сначала ищите нули, единицы и небольшие координаты.\\
Параметр исчез после подстановки & Возьмите другую точку, которая даёт новое уравнение.\\
Получили отрицательное основание $a$ & Для показательной функции такое основание не подходит.\\
Потеряли скобки в $\left(\frac12\right)^{x-3}$ & Подставляемое составное выражение всегда заключайте в скобки.\\
Нашли параметры и остановились & Соберите функцию заново и выполните финальное требование задачи.\\
\bottomrule
\end{tabularx}

\clearpage
\section*{Тренировочный блок — Путь А}
\sectionline
\textbf{10 задач.} Решайте по одному алгоритму: выберите/используйте узловые точки, восстановите параметры, затем выполните требуемое действие. Решения здесь не записаны — только условия.

\begin{enumerate}[label=\textbf{\arabic*.}]
  \item График функции $f(x)=a^x+b$ проходит через точки $(0,-2)$ и $(1,-1)$. Найдите $f(5)$.

  \item График функции $f(x)=a^x+b$ проходит через точки $(0,4)$ и $(1,7)$. Найдите $f(2)$.

  \item График функции $f(x)=a^x+b$ проходит через точки $(0,-3)$ и $(-2,0)$. Найдите $f(3)$.

  \item График функции $f(x)=a^x+b$ проходит через точки $(0,2)$ и $(2,10)$. Решите уравнение $f(x)=28$.

  \item График функции $f(x)=a^x+b$ проходит через точки $(0,-1)$ и $(2,3)$. Найдите $f(4)$.

  \item График функции $f(x)=a^{x+b}$ проходит через точки $(-2,1)$ и $(0,4)$. Найдите $f(3)$.

  \item График функции $f(x)=a^{x+b}$ проходит через точки $(3,1)$ и $(0,8)$. Найдите $f(1)$.

  \item Для функции $f(x)=a^x$ на графике отмечены точки $(0,1)$ и $(-2,9)$. Определите $a$ по информативной точке и найдите $f(3)$.

  \item График функции $f(x)=a^{x+b}$ проходит через точки $(-3,1)$ и $(-1,2)$. Решите уравнение $f(x)=8$.

  \item График функции $f(x)=a^x+b$ проходит через точки $(0,5)$ и $(-1,6)$. Решите уравнение $f(x)=36$.
\end{enumerate}

\vfill
\textit{Самопроверка перед сдачей:} подставьте обе исходные точки в найденную формулу. Если хотя бы одна не подходит, ошибка находится раньше финального вычисления.

\clearpage
\section*{Ответы}
\sectionline
\renewcommand{\arraystretch}{1.35}
\begin{tabularx}{\linewidth}{@{}>{\bfseries}c X@{}}
\toprule
№ & Ответ\\
\midrule
1 & $29$\\
2 & $19$\\
3 & $-\dfrac{31}{8}$\\
4 & $x=3$\\
5 & $23$\\
6 & $32$\\
7 & $4$\\
8 & $a=\dfrac13$, \quad $f(3)=\dfrac1{27}$\\
9 & $x=3$\\
10 & $x=-5$\\
\bottomrule
\end{tabularx}

\vspace{12mm}
\textbf{Ключ к проверке параметров}
\begin{itemize}
  \item В задачах 1--5 и 10 сначала используйте точку с $x=0$: она сразу даёт $1+b$.
  \item В задачах 6, 7 и 9 используйте точку с $y=1$: при $a>0$, $a\ne1$ это означает нулевой показатель степени.
  \item В задаче 8 точка $(0,1)$ не определяет $a$; нужна точка $(-2,9)$.
\end{itemize}

\vfill
\begin{center}
\small\color{softgray}Лёвин Артём Александрович · эксклюзивно для Даниила Кичука · 02.10.26
\end{center}

\end{document}
